\chapter{Related Work}
\section{LLMs and Mathematical Modeling}
\subsection{End-to-End Modeler and Solver Archirecture}
Extensive research has been conducted at the intersection of Large Language Models (LLMs) and mathematical optimization. A prominent example is \textbf{OptLLM} \cite{zhang2024solvinggeneralnaturallanguagedescriptionoptimization}, a framework designed to bridge natural language descriptions and specialized solvers. OptLLM utilizes a modular architecture consisting of three primary components: the \emph{Interaction Refinement Module}, which processes user queries and handles multi-round clarifications; the \emph{Converter Module}, which includes a formulator to translate text into mathematical structures and a coder to generate executable code; and finally, the \emph{Responser Module}, which executes the model via external solvers and interprets the results for the user. This architecture demonstrates the feasibility of integrating disparate modules to transform an LLM into a comprehensive end-to-end optimization system.


\subsection{Natural Langauge to Formal Specification}
Another significant contribution is \textbf{LLMOPT} \cite{jiang2025llmoptlearningdefinesolve}, a learning-based framework that focuses on boosting the generalization of LLMs across various optimization types. LLMOPT introduces a universal \emph{Five-Element Formulation} that reduces complex problem descriptions into five core components: sets, parameters, variables, objectives, and constraints. This structured representation allows the system to generate consistent solver code across diverse domains, including linear, nonlinear, and combinatorial optimization.

\subsection{LLM Benchmark for Optimization Modeling}
While the aforementioned frameworks focus on generation, \textbf{Mamo} \cite{huang2025llmsmathematicalmodelingbridging} provides a specialized benchmark to evaluate these modeling capabilities. Mamo consists of 1,209 meticulously curated questions spanning Ordinary Differential Equations (ODEs), Linear Programming (LP), and Mixed-Integer Linear Programming (MILP). Unlike traditional result-oriented metrics, Mamo utilizes a ``model generation and numerical validation'' approach. An LLM is tasked with formulating a model from the dataset, which is then executed by a solver; the resulting numerical answer is compared against the provided ground truth. This benchmark is particularly valuable for the evaluation phase of this thesis, as it provides a standardized methodology for measuring the accuracy of AI-generated mathematical models.


\section{AI Code Generation}

\subsection{Current State}
Since the introduction of the Transformer architecture \cite{vaswani2023attentionneed}, there has been a paradigm shift in the field of automated software engineering. While Transformers were initially celebrated for their success in Natural Language Processing (NLP), their ability to model long-range dependencies and complex structural patterns has led to significant advancements in Large Language Models (LLMs) for code generation.

Extensive research has explored the transition from general-purpose text generation to domain-specific code synthesis. As detailed in the comprehensive survey by Jiang et al. \cite{Jiang_2026}, modern frameworks leverage massive datasets of source code to perform tasks such as code completion, bug detection, and natural-language-to-code translation. The survey categorizes state-of-the-art models—including GPT-4, CodeLlama, and StarCoder—and examines the underlying training objectives, such as ``Fill-in-the-Middle'' (FIM) and instruction tuning, which are essential for understanding the non-linear structure of programming languages.

\subsection{LLM Benchmarking for Code Generation}
As Large Language Models (LLMs) continue to evolve, the demand for more rigorous evaluation methodologies has led to the development of advanced benchmarking frameworks. One such framework is \textbf{EvalPlus} \cite{liu2023codegeneratedchatgptreally}, which addresses the limitations of traditional code generation benchmarks. 

EvalPlus builds upon established datasets, most notably \emph{HumanEval} \cite{chen2021evaluatinglargelanguagemodels} and \emph{MBPP}. While these original benchmarks provided a foundational method for measuring functional correctness, their test suites are often sparse and fail to catch edge-case failures, allowing incorrect code to pass as correct. EvalPlus mitigates this by automating the synthesis of thousands of additional test cases for each problem. By significantly increasing the test-case density—often by a factor of 80x compared to the original suites—EvalPlus creates a more robust evaluation environment that exposes ``false positives'' in LLM-generated code. 


\section{Human-in-the-loop Architecture}

\textbf{Magentic-UI} \cite{mozannar2025magenticuihumanintheloopagenticsystems} is a research prototype developed by Microsoft Research that introduces a human-centered AI agent designed to execute complex web-based tasks. Unlike fully autonomous systems, Magentic-UI emphasizes a collaborative relationship between the user and the agent. The system receives natural language queries and interprets user intent to formulate a multi-step plan. 



A key contribution of this framework is its implementation of several interaction mechanisms, most notably \emph{co-planning} and \emph{action guards}. During co-planning, the agent reveals its intended steps to the user, who can then audit, alter, or pause the process before execution begins. Furthermore, all critical or irreversible actions—such as submitting sensitive data or performing online transactions—require explicit human approval via action guards. This architecture ensures that operations are conducted within a safe, transparent environment where the user maintains final authority. By prioritizing human oversight and iterative feedback, Magentic-UI serves as a state-of-the-art example of a well-formed human-in-the-loop (HITL) architecture for web-agentic systems.

% \section{Schema Validation}